% ==============================================================================
% WEEK 2
% ==============================================================================
\section{Week 2 -- 23.02.26 Cauchy's Riemann and Differentiation}

\textbf{Recall:} A function $f: \mathbb{C} \rightarrow \mathbb{C}$ is complex differentiable at $z_{0}$ if the limit
$$ f'(z_{0}) = \lim_{z \rightarrow z_{0}} \frac{f(z) - f(z_{0})}{z - z_{0}} $$
exists and is finite. Analogous to the definition of the derivative of a real function in Analysis I \& II.

\subsection{Cauchy-Riemann equations}

\begin{theorem}[Theorem: Cauchy-Riemann equations]
Let $W$ be an open set and $f: W \rightarrow \mathbb{C}$, $z = x + iy \mapsto f(x+iy) = u(x,y) + iv(x,y)$ be a holomorphic function. Then, $u, v \in C^{1}(W)$ (differentiable with continuous derivatives) and:
$$ \frac{\partial u}{\partial x} = \frac{\partial v}{\partial y} $$
$$ \frac{\partial u}{\partial y} = -\frac{\partial v}{\partial x} $$

Conversely, if these are satisfied (and $u, v \in C^1$), then $f$ is holomorphic.
\end{theorem}
We will prove the direct implication (the converse is more difficult).

\textbf{Proof:} If $f$ is holomorphic on $W$, then for $z_0 \in W$:
$$ f'(z_{0}) = \lim_{z \rightarrow z_{0}} \frac{f(z) - f(z_{0})}{z - z_{0}} $$
We will calculate the limit in two ways: $z = x + iy$, $z_{0} = x_{0} + iy_{0}$.

\textbf{1. Approaching $z_{0}$ from the x-direction ($y = y_0 \implies z = x + iy_{0}$):}
\begin{align*}
f'(z_{0}) &= \lim_{x \rightarrow x_{0}} \frac{f(x + iy_{0}) - f(x_{0} + iy_{0})}{x + iy_{0} - (x_{0} + iy_{0})} \\
&= \lim_{x \rightarrow x_{0}} \frac{u(x, y_{0}) + iv(x, y_{0}) - u(x_{0}, y_{0}) - iv(x_{0}, y_{0})}{x - x_{0}} \\
&= \lim_{x \rightarrow x_{0}} \frac{u(x, y_{0}) - u(x_{0}, y_{0})}{x - x_{0}} + i \lim_{x \rightarrow x_{0}} \frac{v(x, y_{0}) - v(x_{0}, y_{0})}{x - x_{0}} \\
&= \frac{\partial u}{\partial x} + i \frac{\partial v}{\partial x} \quad \text{--- (1)}
\end{align*}

\textbf{2. Approaching $z_{0}$ from the y-direction ($x = x_0 \implies z = x_{0} + iy$):}
\begin{align*}
f'(z_{0}) &= \lim_{y \rightarrow y_{0}} \frac{f(x_{0} + iy) - f(x_{0} + iy_{0})}{x_{0} + iy - (x_{0} + iy_{0})} \\
&= \lim_{y \rightarrow y_{0}} \frac{u(x_{0}, y) + iv(x_{0}, y) - u(x_{0}, y_{0}) - iv(x_{0}, y_{0})}{i(y - y_{0})} \\
&= \frac{1}{i} \lim_{y \rightarrow y_{0}} \left( \frac{u(x_{0}, y) - u(x_{0}, y_{0})}{y - y_{0}} + i \frac{v(x_{0}, y) - v(x_{0}, y_{0})}{y - y_{0}} \right) \\
&= \frac{1}{i} \left( \frac{\partial u}{\partial y} + i \frac{\partial v}{\partial y} \right) \\
&= \frac{\partial v}{\partial y} - i \frac{\partial u}{\partial y} \quad \text{--- (2)}
\end{align*}

Since $(1) = (2)$, we have:
$$ \frac{\partial u}{\partial x} = \frac{\partial v}{\partial y} \quad \text{and} \quad \frac{\partial v}{\partial x} = -\frac{\partial u}{\partial y} $$

Also from this proof, expressions for $f'(z_{0})$:
$$ f' = \frac{\partial u}{\partial x} + i\frac{\partial v}{\partial x} = \frac{\partial v}{\partial y} - i\frac{\partial u}{\partial y} $$

\begin{example}[Examples: Cauchy-Riemann]
\textbf{Ex 1.} $f(z) = z^{2} = (x+iy)^{2} = x^{2} - y^{2} + i(2xy)$
$u(x,y) = x^{2} - y^{2}$, $v(x,y) = 2xy$.
Is it holomorphic over $\mathbb{C}$? We check the Cauchy-Riemann eqns:
\begin{align*}
\frac{\partial u}{\partial x} &= 2x, \quad \frac{\partial v}{\partial x} = 2y \\
\frac{\partial u}{\partial y} &= -2y, \quad \frac{\partial v}{\partial y} = 2x
\end{align*}
C-R satisfied! And $f'(z) = \frac{\partial u}{\partial x} + i\frac{\partial v}{\partial x} = 2x + i(2y) = 2z$.

\textbf{Ex 2.} The complex conjugate: $f(z) = \overline{z} = x - iy$.
$u = x, v = -y$.
\begin{align*}
\frac{\partial u}{\partial x} &= 1, \quad \frac{\partial v}{\partial x} = 0 \\
\frac{\partial u}{\partial y} &= 0, \quad \frac{\partial v}{\partial y} = -1
\end{align*}
Here $\partial_x u \neq \partial_y v$. Not holomorphic!
Another way of seeing the same fact: If we try to compute the limit $\lim_{z \rightarrow z_{0}} \frac{\overline{z} - \overline{z_{0}}}{z - z_{0}}$:
\begin{itemize}
    \item For $z = x + iy_{0}$, $z_{0} = x_{0} + iy_{0}$: $\lim_{x \to x_{0}} \frac{x - x_{0}}{x - x_{0}} = 1$.
    \item For $z = x_{0} + iy$, $z_{0} = x_{0} + iy_{0}$: $\lim_{y \to y_{0}} \frac{-iy + iy_{0}}{iy - iy_{0}} = -1$.
\end{itemize}

\textbf{Ex 3.} $f: \mathbb{C} \setminus \{0\} \to \mathbb{C}$, $f(z) = \frac{1}{z}$.
We have: $f(x+iy) = \frac{\overline{z}}{|z|^{2}} = \frac{x}{x^{2}+y^{2}} - i\frac{y}{x^{2}+y^{2}}$.
So:
\begin{align*}
\frac{\partial u}{\partial x} &= \frac{x^{2}+y^{2} - 2x^{2}}{(x^{2}+y^{2})^{2}} = \frac{y^2-x^2}{(x^2+y^2)^2} \\
\frac{\partial v}{\partial x} &= \frac{2xy}{(x^{2}+y^{2})^{2}} \\
\frac{\partial u}{\partial y} &= -\frac{2xy}{(x^{2}+y^{2})^{2}} \\
\frac{\partial v}{\partial y} &= \frac{-(x^{2}+y^{2}) + 2y^{2}}{(x^{2}+y^{2})^{2}} = \frac{y^2-x^2}{(x^2+y^2)^2}
\end{align*}
So C-R eqns are satisfied.
And $f' = \frac{\partial u}{\partial x} + i\frac{\partial v}{\partial x} = \frac{y^{2}-x^{2}}{(x^{2}+y^{2})^{2}} + i\frac{2xy}{(x^{2}+y^{2})^{2}}$.
This is $= -\frac{1}{z^{2}}$. Indeed: $-\frac{1}{z^{2}} = -\frac{\overline{z}^{2}}{|z|^{4}} = -\frac{(x-iy)^{2}}{(x^{2}+y^{2})^{2}} = \dots$

\textbf{Ex 4.} $f(z) = e^{z} = e^{x+iy} = e^{x}\cos(y) + ie^{x}\sin(y)$.
So:
\begin{align*}
\frac{\partial u}{\partial x} &= e^{x}\cos(y), \quad \frac{\partial v}{\partial x} = e^{x}\sin(y) \\
\frac{\partial u}{\partial y} &= -e^{x}\sin(y), \quad \frac{\partial v}{\partial y} = e^{x}\cos(y)
\end{align*}
So C-R eqns are satisfied.
And $f' = \frac{\partial u}{\partial x} + i\frac{\partial v}{\partial x} = e^{x}\cos(y) + ie^{x}\sin(y) = e^{z}$.
\end{example}

\subsection{Rules of differentiation for complex derivatives}

Let $f, g: W \rightarrow \mathbb{C}$ be holomorphic. Then:
\begin{itemize}
    \item $(f+g)' = f' + g'$
    \item $(f \cdot g)' = f' \cdot g + f \cdot g'$
    \item $\left(\frac{f}{g}\right)' = \frac{f' \cdot g - f \cdot g'}{g^{2}}$ (when $g \neq 0$)
    \item If the image of $g$ is inside $W$: $(f \circ g)'(z) = f'(g(z)) \cdot g'(z)$.
\end{itemize}

These relations can be proved in exactly the same way as in the real case. E.g.
\begin{align*}
(f \cdot g)'(z_{0}) &= \lim_{z \to z_{0}} \frac{f(z) \cdot g(z) - f(z_{0})g(z_{0})}{z - z_{0}} \\
&= \lim_{z \to z_{0}} \frac{f(z)g(z) - f(z_{0})g(z) + f(z_{0})g(z) - f(z_{0})g(z_{0})}{z - z_{0}} \\
&= \lim_{z \to z_{0}} \left( \frac{f(z)-f(z_{0})}{z - z_{0}}g(z) + f(z_{0})\frac{g(z)-g(z_{0})}{z - z_{0}} \right) \\
&= f'(z_{0})g(z_{0}) + f(z_{0})g'(z_{0})
\end{align*}

With these rules and the examples we computed earlier, it follows that the usual formulas for polynomials, trigonometric functions, etc. hold.
E.g. $(z^{k})' = k \cdot z^{k-1}$ ($k \in \mathbb{Z}$), $(\cos(z))' = -\sin(z)$.

\subsection{Aside: A geometric property of holomorphic functions}

If $f: W \rightarrow \mathbb{C}$ is holomorphic, when viewed as a map $f: \mathbb{R}^{2} \rightarrow \mathbb{R}^{2}$, the Jacobian matrix
$$ Df = \begin{pmatrix} \partial_{x}u & \partial_{y}u \\ \partial_{x}v & \partial_{y}v \end{pmatrix} $$
tells us how $f$ acts on infinitesimal rectangles.
By the Cauchy-Riemann equations, $Df$ takes the form of $\lambda \cdot \text{Rotation matrix}$.
So an infinitesimal square is mapped to a square (no uneven stretching, preserving angles -- this is called a \textit{conformal mapping}).

\subsection{The complex logarithm}

\begin{definition}[Definition: Complex Logarithm]
\textbf{Recall:} The principal argument of a complex number $z \neq 0$: $-\pi < \text{Arg}(z) \leq \pi$.
It is continuous except on the negative real axis (where it jumps from $\pi$ to $-\pi$).

How would we define the complex logarithm? We want $e^{\log(z)} = z$.
If $\log(z) = a + ib$, then $z = e^{a+ib}$.
$|z| \cdot e^{i\text{Arg}(z)} = e^{a} \cdot e^{ib}$.
It makes sense to set $|z| = e^{a}$ and $\text{Arg}(z) = b$ (up to $2\pi k$).

So: We define $\log(z) = \log|z| + i\text{Arg}(z)$.
\begin{itemize}
    \item $\log(z)$ is not defined for $z=0$.
    \item $\log(z)$ is discontinuous on the negative real axis.
    \item But: $\log(z)$ is continuous on $\mathbb{C} \setminus \{z \in \mathbb{C} : \text{Im}(z) = 0, \text{Re}(z) \leq 0\}$. This excluded line is the \textit{Branch cut}. (Had we defined $\text{Arg}(z)$ differently, we would have a different branch cut).
\end{itemize}
\end{definition}

Real and imaginary parts: $\log(z) = \log|z| + i\text{Arg}(z)$.
If $z = x+iy$: If $x > 0$,
$$ \log(z) = \log\sqrt{x^{2}+y^{2}} + i\arctan\left(\frac{y}{x}\right) $$
where $u(x,y) = \log\sqrt{x^2+y^2}$ and $v(x,y) = \arctan(y/x)$.

Easy to verify: $\log(z)$ is holomorphic on $\mathbb{C} \setminus \{z : \text{Im}(z) = 0, \text{Re}(z) \leq 0\}$ with $(\log z)' = \frac{1}{z}$.

\textit{Remark:} For positive real numbers, $\log(z_{1} \cdot z_{2}) = \log(z_{1}) + \log(z_{2})$.
For general complex numbers (due to discontinuity), this is \textbf{no longer true}.
If $z_{1} = z_{2} = -1$, $\log(z_{1}) = \log(z_{2}) = i\pi$, but $\log(z_{1} \cdot z_{2}) = \log(1) = 0$.
But: It is true up to $2\pi i$.

\subsection{Power function with complex exponent}

Let $\gamma \in \mathbb{C}$. Define:
$$ f(z) = z^{\gamma} = e^{\log(z^{\gamma})} = e^{\gamma \cdot \log(z)} $$
$= e^{\gamma \cdot \log|z| + i\gamma \cdot \text{Arg}(z)}$.

It is holomorphic over $\mathbb{C} \setminus \{z : \text{Im}(z) = 0, \text{Re}(z) \leq 0\}$ (as the composition of $e^{\gamma z}$ with $\log z$).
When $\gamma \in \mathbb{Z}$: it is continuous (in fact, holomorphic) for all $z \neq 0$.